\documentclass[../../main.tex]{subfiles}
\graphicspath{{../../image/}{../../figure/}}
\begin{document}

\section{注记}\label{section:ch13-注记}

随机选择法由 Glimm 的奠基性论文 Glimm [1] 提出。正是在该工作中，相容性（\cref{section:ch13-紧性与相容性}）、波相互作用（\cref{section:ch13-真非线性波相互作用}）与 Glimm 泛函（\cref{section:ch13-glimm泛函}）的思想被首次引入，而\cref{theorem:ch13-小总变差存在性}也是对真非线性系统首先建立的。如下一章所见，正是 Glimm 型泛函为其他解逼近格式的紧性提供了关键估计。此外，即使在一般 $BV$ 解的框架中，Glimm 泛函也可定义并卓有成效地使用；见\cref{chapter:ch14-前沿追踪与标准riemann半群}。

对等温气体动力学特殊系统\eqref{align:ch13-isothermal}，大总变差解的构造归功于 Nishida [1]。它基于识别一个 Riemann 不变量的时间递减泛函，从而给出解的变差界。关于大初始数据（或至少中等偏大初始数据）解的相关构造，见 Bakhvarov [1]、DiPerna [1,2]、Nishida 与 Smoller [1]、Cai Zhong Li [1]、Luskin 与 Temple [1]、Serre [11]、Ying 与 Wang [1]、Amadori 与 Guerra [2]、Asakura 与 Corli [1]，以及 Holden、Risebro 与 Sande [1]。另一方面，Chen 与 Jenssen [1] 表明，具备此种结构的系统类是贫乏的。

波对消（引入于\cref{section:ch13-真非线性波相互作用}）对解的存在性与长时间行为的影响，首先在 Glimm 与 Lax [1] 的开创性论文集中得到展示，该文献已在\cref{chapter:ch12-两方程真非线性组}中引用。

\cref{section:ch13-上确界界}中概述的上确界界的推导取自 R. Young [1] 的学位论文，读者可在其中找到技术细节。事实上，该工作为 Cauchy 问题引入了一个新的长度尺度，在特定情形下可用于放宽对多于两个守恒律的某些系统的初始数据小总变差要求。沿此方向，见 Temple [7]、Temple 与 Young [1,2]，以及 Cheverry [3]。局部或整体解在大总变差初始数据下的构造，亦见于 Alber [1] 与 Schochet [3,4]。

Glimm 泛函被 Tai-Ping Liu [15] 推广到非真非线性系统，他首先认识到同一特征族趋近波之间夹角所起的重要作用。\cref{section:ch13-一般系统}中给出的概述遵循 Iguchi 与 LeFloch [1] 以及 Tai-Ping Liu 与 Tong Yang [6] 的较近期工作。

波划分方法在 Tai-Ping Liu [7] 中为真非线性系统发展，并在 Tai-Ping Liu [15] 中用于建立算法对等分布序列的确定性相容性（\cref{theorem:ch13-等分布相容}）。进而，跟随\cref{chapter:ch12-两方程真非线性组}中引用的 DiPerna [3] 关于真非线性系统的工作，

关于随机选择格式收敛速度的早期工作，见 Hoff 与 Smoller [1]、Tong Yang [3]、Ye 与 Lin [1]，以及 Hua 与 Yang [1]。对``优等分布''序列，$L^1$ 中更精细的收敛速度 $o(h^{1/2}|\log h|)$ 由 Bressan 与 Marson [3] 对真非线性系统建立。

Hua、Jiang 与 Yang [1]、Hua 与 Yang [2]、Ancona 与 Marson [8,9]、Bianchini 与 Modena [1,2,3]、Modena 与 Bianchini [1] 以及 Modena [1] 的较近期工作，基于``二次 Glimm 泛函''，为探测格式的收敛速度及所得解的局部结构提供了非常有力的工具。

另一种 Glimm 泛函由 Caravenna [1] 提出，基于极大熵产生，至少对标量守恒律有效。

当 $t\to\infty$，\eqref{align:ch13-ivp}的解趋近以\cref{chapter:ch09-容许波扇与riemann问题}为数据的 Riemann 问题的解，
其中 $U_L=U_0(-\infty)$，$U_R=U_0(+\infty)$；参见 Tai-Ping Liu [9,11,15]。较近的阐述见 Tai-Ping Liu [28] 以及 Tai-Ping Liu 与 Tong Yang [6]。

关于随机选择法的推广与应用，文献浩繁。混合型系统见 Pego 与 Serre [1]、LeFloch [3] 以及 Corli 与 Sablé-Tougeron [3]。初边值问题参见 Tai-Ping Liu [11]、Luskin 与 Temple [1]、Nishida 与 Smoller [2]、Dubroca 与 Gallice [1]、Sablé-Tougeron [1] 以及 Frid [1]。含强激波的解见 Sablé-Tougeron [2]、Corli 与 Sablé-Tougeron [1,2]、Asakura [1]、Corli [2] 以及 Schochet [3,4]。周期解见 Frid [4] 以及 Frid 与 Perepelitsa [1]。在气体动力学中的应用见 Tai-Ping Liu [4,5,12,16,17]、Temple [1] 以及 Tong Yang [2]。气体动力学中真空的效应见 Liu 与 Smoller [1] 以及 Long-Wei Lin [1,2]。在相对论理论中的应用见 Barnes、LeFloch、Schmidt 与 Stewart [1]。非守恒形式的系统见 LeFloch [2] 以及 LeFloch 与 Liu [1]。弱 $L^p$ 稳定性由 Temple [6] 建立。更多文献见 Smoller [3]、Serre [11] 与 LeFloch [5] 的专著。

\end{document}
